%==============================================================
%  lecons/analyse/series_asymptotiques.tex — Séries asymptotiques
%==============================================================

\lecontitle{Séries asymptotiques}{Analyse réelle — Développements de Poincaré et troncature optimale}

%=== NOTATION LOCALE =========================================
% (aucune : \sim, \mathcal{O}, o sont fournis par amsmath/amssymb)

%=============================================================
%  § 1 — DÉFINITION ET RÉSULTATS FONDAMENTAUX
%=============================================================
\secnum{navyblue}{1}{Définition et résultats fondamentaux}

\begin{tcolorbox}[thmbox]
  \textbf{Définition (développement asymptotique de Poincaré).}%
  \enspace\index{série asymptotique}\index{développement asymptotique}
  Soit $f$ définie au voisinage de $+\infty$ et $(a_n)_{n\geq 0}$ une suite de
  scalaires. On dit que la série $\sum a_n x^{-n}$ est un \emph{développement
  asymptotique} de $f$ en $+\infty$, et l'on écrit
  \[
    f(x) \;\sim\; \sum_{n=0}^{+\infty} \frac{a_n}{x^{n}}
    \qquad (x\to +\infty),
  \]
  si pour tout entier $N\geq 0$ le reste vérifie
  \[
    f(x) - \sum_{n=0}^{N} \frac{a_n}{x^{n}}
    \;=\; \mathcal{O}\!\left(\frac{1}{x^{N+1}}\right)
    \qquad (x\to +\infty),
  \]
  ce qui équivaut à $f(x)-\sum_{n=0}^{N}a_n x^{-n}=o(x^{-N})$.

  \medskip
  \textbf{Point capital.} La définition ne suppose \emph{aucune} convergence de
  la série $\sum a_n x^{-n}$ : elle fixe $N$ et fait tendre $x\to+\infty$, à
  l'inverse d'une série convergente qui fixe $x$ et fait $N\to+\infty$. Une série
  asymptotique peut diverger pour \emph{tout} $x$.

  \medskip
  \textbf{Théorème (propriétés fondamentales).}
  \begin{itemize}
    \item \textit{Unicité des coefficients.} Si $f$ admet un développement
    asymptotique, ses coefficients sont uniques, donnés par
    \[
      a_0=\lim_{x\to+\infty} f(x),\qquad
      a_n=\lim_{x\to+\infty} x^{n}\Bigl(f(x)-\sum_{k=0}^{n-1}a_k x^{-k}\Bigr).
    \]
    \item \textit{Non-unicité de la fonction.} Deux fonctions distinctes peuvent
    partager le même développement : $f$ et $f+e^{-x}$ ont le même, car
    $e^{-x}=o(x^{-N})$ pour tout $N$ (termes « au-delà de tout ordre »).
    \item \textit{Linéarité et produit.} L'ensemble des fonctions admettant un
    développement asymptotique est stable par combinaison linéaire et par
    produit, avec développement obtenu par les opérations formelles
    correspondantes.
    \item \textit{Intégration.} On peut intégrer terme à terme un développement
    asymptotique ; on ne peut pas le \emph{dériver} sans hypothèse
    supplémentaire.
    \item \textit{Troncature optimale.} Pour une série asymptotique divergente,
    l'erreur de la somme partielle $S_N(x)$ est de l'ordre du premier terme
    négligé ; elle décroît puis croît avec $N$. Tronquer au plus petit terme
    minimise l'erreur (\emph{superasymptotique}).
  \end{itemize}
\end{tcolorbox}

\begin{significationintuitive}
Convergence et asymptotique échangent l'ordre de deux limites. Une série
convergente promet : « fixe la précision voulue, je l'atteins en prenant assez
de termes ». Une série asymptotique promet l'inverse : « fixe le nombre de
termes, l'erreur devient minuscule quand $x$ grandit ». C'est pourquoi quelques
termes d'une série \emph{divergente} suffisent souvent à calculer une fonction
avec une précision spectaculaire — à condition de s'arrêter au bon moment, près
du plus petit terme.
\end{significationintuitive}

\decsep{}

%=============================================================
%  § 2 — DÉMONSTRATION
%=============================================================
\secnum{navyblue}{2}{Démonstration}

\begin{ideepreuve}
On construit l'exemple canonique d'Euler $F(x)=\int_0^{+\infty}e^{-xt}(1+t)^{-1}\,
\mathrm{d}t$. Le développement géométrique de $1/(1+t)$ avec reste intégral, suivi
de $\int_0^\infty e^{-xt}t^k\,\mathrm{d}t=k!/x^{k+1}$, fournit la série
$\sum(-1)^k k!\,x^{-k-1}$ et un reste majoré par le terme suivant. La série
diverge pour tout $x$, mais le contrôle du reste donne l'asymptotique et la
troncature optimale.
\end{ideepreuve}

\vspace{10pt}
\rubriq{Preuve}

\begin{mdframed}[style=proofbar]
  \textit{Étape 1 — Unicité des coefficients.}
  Supposons $f(x)\sim\sum a_n x^{-n}$. En faisant $x\to+\infty$ dans
  $f(x)-\sum_{k<n}a_k x^{-k}=a_n x^{-n}+o(x^{-n})$, on obtient
  $x^{n}\bigl(f(x)-\sum_{k<n}a_k x^{-k}\bigr)\to a_n$. Les $a_n$ se calculent
  donc de proche en proche à partir de $f$ seule : ils sont uniques.

  \medskip
  \textit{Étape 2 — Construction de l'exemple d'Euler.}
  Pour $x>0$ posons
  \[
    F(x)=\int_{0}^{+\infty}\frac{e^{-xt}}{1+t}\,\mathrm{d}t .
  \]
  L'intégrale converge (domination par $e^{-xt}$). De l'identité, pour tout
  $N\geq 0$,
  \[
    \frac{1}{1+t}=\sum_{k=0}^{N-1}(-1)^k t^{k}
      +\frac{(-1)^{N}t^{N}}{1+t},
  \]
  et de $\displaystyle\int_0^{+\infty}e^{-xt}t^{k}\,\mathrm{d}t=\frac{k!}{x^{k+1}}$,
  on tire
  \[
    F(x)=\sum_{k=0}^{N-1}\frac{(-1)^{k}\,k!}{x^{k+1}}
      \;+\; R_N(x),
    \qquad
    R_N(x)=(-1)^{N}\!\int_{0}^{+\infty}\frac{e^{-xt}t^{N}}{1+t}\,\mathrm{d}t .
  \]

  \medskip
  \textit{Étape 3 — Majoration du reste et caractère asymptotique.}
  Comme $0\leq \dfrac{1}{1+t}\leq 1$ sur $[0,+\infty[$,
  \[
    |R_N(x)|\;\leq\;\int_{0}^{+\infty}e^{-xt}t^{N}\,\mathrm{d}t
    \;=\;\frac{N!}{x^{N+1}} .
  \]
  Le reste est donc majoré, en valeur absolue, par le \emph{premier terme
  négligé}. En particulier $R_N(x)=\mathcal{O}(x^{-N-1})$, ce qui établit
  \[
    F(x)\;\sim\;\sum_{k=0}^{+\infty}\frac{(-1)^{k}\,k!}{x^{k+1}}
    \qquad (x\to+\infty).
  \]

  \medskip
  \textit{Étape 4 — Divergence de la série.}
  Pour $x>0$ fixé, le terme général $u_k=k!/x^{k+1}$ vérifie
  $u_{k+1}/u_k=(k+1)/x\to+\infty$ : il ne tend pas vers $0$, donc la série
  $\sum(-1)^k k!\,x^{-k-1}$ \emph{diverge grossièrement} pour tout $x$.
  Convergence et asymptotique sont bien deux notions indépendantes.

  \medskip
  \textit{Étape 5 — Troncature optimale.}
  L'erreur commise en s'arrêtant après $N$ termes est de l'ordre de
  $u_N=N!/x^{N+1}$. Le rapport $u_{N+1}/u_N=(N+1)/x$ est $<1$ tant que $N<x-1$ et
  $>1$ ensuite : les termes décroissent puis croissent, le plus petit étant
  atteint vers $N\approx x$. Par la formule de Stirling,
  $N!/x^{N+1}\approx \sqrt{2\pi x}\,e^{-x}/x$ pour $N\approx x$, d'où une erreur
  optimale de l'ordre de $e^{-x}$ : minuscule, bien que la série diverge.
  \hfill$\blacksquare$
\end{mdframed}

\decsep{}

%=============================================================
%  § 3 — EXEMPLES, CONTRE-EXEMPLES ET EXERCICES
%=============================================================
\secnum{exgreen}{3}{Exemples, contre-exemples et exercices}

\rubriq{Exemples}

\begin{mdframed}[style=exbar]
  \textbf{1. Série d'Euler.}\enspace\index{série d'Euler}%
  L'exemple fondateur, dû à Euler (1760), est
  \[
    F(x)=\int_{0}^{+\infty}\frac{e^{-xt}}{1+t}\,\mathrm{d}t
    \;\sim\;\frac{1}{x}-\frac{1}{x^{2}}+\frac{2!}{x^{3}}-\frac{3!}{x^{4}}+\cdots
  \]
  La série diverge partout, mais trois ou quatre termes calculent $F(x)$ avec une
  excellente précision dès que $x$ est grand.

  \smallskip\noindent
  \textbf{2. Série de Stirling.}\enspace\index{formule de Stirling}%
  Le développement de la fonction Gamma,
  \[
    \ln\Gamma(x)\;\sim\;\Bigl(x-\tfrac12\Bigr)\ln x-x+\tfrac12\ln(2\pi)
      +\sum_{n\geq 1}\frac{B_{2n}}{2n(2n-1)\,x^{2n-1}},
  \]
  a sa partie en série \emph{divergente} (les nombres de Bernoulli $B_{2n}$
  croissent très vite), mais asymptotique : c'est elle qui rend la formule de
  Stirling si précise pour quelques termes.

  \smallskip\noindent
  \textbf{3. Fonction d'erreur complémentaire.}\enspace%
  Pour $x\to+\infty$,
  \[
    \int_{x}^{+\infty}e^{-t^{2}}\,\mathrm{d}t
    \;\sim\;\frac{e^{-x^{2}}}{2x}\Bigl(1-\frac{1}{2x^{2}}+\frac{3}{4x^{4}}
      -\cdots\Bigr),
  \]
  obtenue par intégrations par parties successives ; là encore la série entre
  parenthèses diverge mais approxime remarquablement la queue gaussienne.
\end{mdframed}

\vspace{6pt}
\rubriq{Contre-exemples et pièges}

\begin{mdframed}[style=cexbar]
  \textbf{Asymptotique $\neq$ convergente.}\enspace%
  On ne peut pas atteindre une précision arbitraire à $x$ fixé en ajoutant des
  termes : passé le plus petit terme, l'erreur \emph{réaugmente}. Sommer
  « tous » les termes n'a aucun sens.

  \smallskip\noindent
  \textbf{La fonction n'est pas déterminée.}\enspace%
  $F(x)$ et $F(x)+e^{-x}$ ont le même développement asymptotique : un
  développement asymptotique ne caractérise pas la fonction (termes
  exponentiellement petits invisibles à tous les ordres).

  \smallskip\noindent
  \textbf{Pas de dérivation terme à terme.}\enspace%
  De $f(x)\sim\sum a_n x^{-n}$ on ne peut pas déduire en général
  $f'(x)\sim -\sum n\,a_n x^{-n-1}$ : ajouter $e^{-x}\sin(e^{x})$, de
  développement asymptotique nul mais dont la dérivée
  $\cos(e^{x})-e^{-x}\sin(e^{x})$ ne tend pas vers $0$, détruit la relation. La
  dérivation exige une hypothèse de régularité supplémentaire.

  \smallskip\noindent
  \textbf{Troncature naïve.}\enspace%
  Prendre $N$ « le plus grand possible » est contre-productif : il faut s'arrêter
  près de $N\approx x$ (plus petit terme), pas au-delà.
\end{mdframed}

\vspace{6pt}
\exolabel

\begin{mdframed}[style=exobar]
  \begin{enumerate}
    \item Démontrer l'unicité des coefficients d'un développement asymptotique et
          calculer $a_0,a_1,a_2$ pour $f(x)=\dfrac{x}{x-1}$ en $+\infty$.

    \item Établir en détail, pour $F(x)=\int_0^{+\infty}\dfrac{e^{-xt}}{1+t}\,
          \mathrm{d}t$, l'égalité
          $F(x)=\sum_{k=0}^{N-1}\dfrac{(-1)^k k!}{x^{k+1}}+R_N(x)$ avec
          $|R_N(x)|\leq N!/x^{N+1}$.

    \item En déduire que la série d'Euler diverge pour tout $x>0$, tout en étant
          asymptotique à $F$.

    \item (Troncature optimale) Montrer que $u_N=N!/x^{N+1}$ est minimal pour
          $N\approx x$ et que la valeur minimale est de l'ordre de $e^{-x}$.
          \textit{Indication :} formule de Stirling.

    \item Obtenir le développement asymptotique de $\displaystyle\int_x^{+\infty}
          e^{-t^2}\,\mathrm{d}t$ par intégrations par parties successives et
          majorer le reste.

    \item (Non-unicité) Vérifier que $x\mapsto e^{-x}$ admet le développement
          asymptotique nul en $+\infty$, et en déduire que deux fonctions de
          même développement diffèrent d'un terme $o(x^{-N})$ pour tout $N$.
  \end{enumerate}
\end{mdframed}

\leconfooter{L.\ Euler (1760) ; H.\ Poincaré, \textit{Sur les intégrales irrégulières} (1886)}
